\section{GraphRNN (20 points)}

Graph generation models produce entire graphs for tasks such as predicting
graph evolution, generating new graph instances, and detecting anomalous
graphs. GraphRNN generates a graph by sequentially adding nodes and edges.
Given a graph \(G\) and a node ordering \(\pi\), it maps the graph to a unique
sequence of node- and edge-addition decisions \(S^\pi\).

When a new node is added, GraphRNN predicts whether it connects to each
existing node. For example, the node-level step for the fourth node consists
of the three binary decisions
\[
  S_{4,1}^{\pi},\qquad S_{4,2}^{\pi},\qquad S_{4,3}^{\pi}.
\]

\subsection{BFS Generation Sequence (12 points)}

Consider the grid graph shown in Question 1.1 of the official homework PDF.
GraphRNN adds its nodes in breadth-first-search order, starting at node \(A\).
When exploring a node's neighbors, visit them in alphabetical order. For each
new node, report a \(1\) for every edge that should be added to an existing
node and a \(0\) otherwise.

Give the node order followed by the edge-level predictions for every node.

\begin{solution}
\textbf{Node order.}
Starting from \(A\), the alphabetically ordered neighbors of \(A\) are
\(B,D\).  Processing \(B\) next discovers \(C,E\), while processing \(D\)
discovers no new node.  Finally, processing \(C\) discovers \(F\).  Therefore
the BFS ordering is
\[
  \pi=(A,B,D,C,E,F).
\]

\textbf{Edge-level predictions.}
For each new node, the entries below are listed against the preceding nodes
in the order \(\pi\):
\[
\begin{array}{c|c|c}
\text{new node} & \text{preceding nodes} & \text{predictions} \\ \hline
A & \text{none} & \text{none} \\
B & A & (1) \\
D & A,B & (1,0) \\
C & A,B,D & (0,1,0) \\
E & A,B,D,C & (0,1,1,0) \\
F & A,B,D,C,E & (0,0,0,1,1)
\end{array}
\]
Here \(1\) denotes an edge to the corresponding preceding node and \(0\)
denotes no edge.
\end{solution}

\subsection{Advantages of BFS Ordering (8 points)}

Explain two advantages of using a BFS node ordering for graph generation
instead of a random node ordering. You may use claims from the original
GraphRNN paper.

\begin{solution}
\begin{enumerate}[label=\textbf{\arabic*.},leftmargin=*]
  \item \textbf{BFS greatly reduces the number of redundant sequential
    representations of the same graph.}

    Let \(G\) have \(n\) nodes and let \(\Pi\) be the set of its \(n!\) node
    permutations.  With an unrestricted ordering, every \(\pi\in\Pi\) can
    produce a different training sequence
    \[
      S^\pi=f_S(G,\pi).
    \]
    These sequences describe the same unlabeled graph, but an autoregressive
    model sees them as different ordered targets.  It would therefore have to
    spread probability mass over, and learn the dependencies in, potentially
    \(n!\) different representations of one graph.

    GraphRNN instead applies a deterministic BFS procedure after choosing a
    starting node and a tie-breaking order:
    \[
      b_G(\pi)=\operatorname{BFS}(G,\pi),\qquad
      S^\pi=f_S\bigl(G,b_G(\pi)\bigr).
    \]
    The map \(b_G\) is many-to-one.  For example, changing the relative
    positions in \(\pi\) of vertices that BFS does not compare can leave the
    starting node, every neighbor-discovery order, and hence the complete BFS
    output unchanged.  Thus, if
    \[
      \mathcal B(G)
      =\{\,\operatorname{BFS}(G,\pi):\pi\in\Pi\,\},
    \]
    then
    \[
      |\mathcal B(G)|\leq |\Pi|=n!,
    \]
    and the inequality is strict for many graphs.  The model only needs to
    learn the sequences induced by \(\mathcal B(G)\), rather than those
    induced by every arbitrary permutation.

    This restriction does not remove any graph from the model's support.
    Every graph has at least one BFS traversal, and the full edge decisions
    associated with any such traversal reconstruct exactly the original
    graph.  BFS therefore removes only redundant orderings, not graph
    structures.  Fewer equivalent targets mean less ordering-induced
    variation, less probability mass split among artificial representations,
    and an easier sequence-learning problem.  The reduction need not be
    strict in the worst case---a star graph can still have \(n!\) BFS
    orderings---but it is substantial for many graph families.

  \item \textbf{BFS localizes possible edges and shortens every edge-generation
    sequence.}

    Under an arbitrary ordering \(v_1,\ldots,v_n\), the adjacency vector for
    the newly added node \(v_i\) is
    \[
      S_i=(A_{1,i},A_{2,i},\ldots,A_{i-1,i})^\top.
    \]
    It has \(i-1\) entries because, without additional structure, \(v_i\)
    might be adjacent to any preceding node.  Consequently, generating an
    \(n\)-node graph requires
    \[
      \sum_{i=2}^{n}(i-1)
      =\frac{n(n-1)}{2}
      =\Theta(n^2)
    \]
    binary edge predictions.

    BFS gives a stronger locality property.  Consider the instant at which a
    new vertex \(v_i\) is first discovered.  Let \(\mathcal F_i\) consist of
    the vertex currently being expanded together with the already discovered
    but not yet expanded vertices in the BFS queue.  We claim that every
    preceding neighbor of \(v_i\) belongs to \(\mathcal F_i\).

    To prove the claim, suppose instead that \(v_i\) were adjacent to a
    preceding vertex \(u\notin\mathcal F_i\).  Such a vertex \(u\) was fully
    processed before \(v_i\) was discovered.  BFS examines every edge incident
    to \(u\) while processing \(u\).  It would therefore have discovered and
    enqueued \(v_i\) at that earlier time, contradicting the assumption that
    \(v_i\) is first discovered now.  Hence old, completed parts of the BFS
    traversal cannot contain a neighbor of the new node; all possible
    preceding neighbors lie in the current frontier.

    Because BFS processes nodes in FIFO discovery order, this frontier is a
    suffix of the preceding node sequence.  Let
    \[
      M=\max_i |\mathcal F_i|
    \]
    be its maximum size.  It is therefore sufficient to retain at most the
    last \(M\) entries of each adjacency vector:
    \[
      S_i=
      \bigl(
        A_{\max(1,i-M),i},\ldots,A_{i-1,i}
      \bigr)^\top,
    \]
    with padding when \(i-1<M\).  All omitted entries are guaranteed to be
    zero by the BFS argument above.  The total number of predictions is now
    bounded by
    \[
      \sum_{i=2}^{n}\min(M,i-1)
      \leq M(n-1)
      =O(Mn).
    \]
    The GraphRNN analysis relates \(M\) to the largest BFS layer:
    \[
      M
      =O\left(
        \max_{1\leq d\leq\operatorname{diam}(G)}
        \bigl|\{v:\operatorname{dist}(v,v_1)=d\}\bigr|
      \right).
    \]
    Thus, whenever the BFS frontier is much smaller than the whole graph
    (\(M\ll n\)), generation is sub-quadratic.  For example, a path has
    \(M=O(1)\), giving \(O(n)\) predictions, while a roughly square grid has
    \(M=O(\sqrt n)\), giving \(O(n^{3/2})\).  Dense graphs or stars may have
    \(M=\Theta(n)\), so the worst case remains quadratic.

    Besides reducing computation and memory, the fixed and shorter
    edge-level sequences remove many guaranteed-zero decisions and shorten
    the dependencies that the edge-level RNN must learn.  BFS therefore makes
    both training and sampling more scalable than a random ordering, which
    can place adjacent vertices arbitrarily far apart.
\end{enumerate}
\end{solution}
